\documentclass[12pt]{article}
\usepackage[T1]{fontenc}
\usepackage[utf8]{inputenc}
\usepackage{lmodern}
\usepackage{textcomp}
\usepackage{enumitem}
\usepackage{geometry}
\geometry{margin=1in}
\begin{document}
\begin{center}
Answer Key - Business Administration - Undergraduate Level Assessment 25 \
\small (Answers are ordered exactly as in the question paper.)
\end{center}

\section*{Multiple Choice}
\begin{enumerate}[label=\arabic*.]
\item \textbf{Answer:} B
\\ \textit{Options:} A) Non-publics; B) Latent publics; C) Inactive publics; D) Complacent publics
\\ \textit{Note:} The term "latent publics" refers to stakeholders who have a potential interest in an issue related to an organization but are currently unaware of their connection or the implications of that issue, distinguishing them from other types of publics who are either aware or actively engaged.
\item \textbf{Answer:} B
\\ \textit{Options:} A) Shannon and Weaver model; B) Osgood and Schramm model; C) Westley and MacLean model; D) Maletzke model
\\ \textit{Note:} The Osgood and Schramm model is correct because it emphasizes the interactive process of encoding and decoding messages between the sender and receiver, highlighting the roles of perception and interpretation in communication.
\item \textbf{Answer:} B
\\ \textit{Options:} A) Political, Legal, Ethical and Philanthropic; B) Economic, Legal, Ethical and Philanthropic; C) Economic, Legal, Cultural and Philanthropic; D) Economic, Legal, Ethical and Environmental
\\ \textit{Note:} The correct answer reflects Carroll's Pyramid of Corporate Social Responsibility, which categorizes CSR into four levels: economic responsibilities as the foundation, followed by legal, ethical, and philanthropic responsibilities that organizations should strive to fulfill.
\item \textbf{Answer:} B
\\ \textit{Options:} A) Team culture; B) Collaborative culture; C) Group culture; D) Collective culture
\\ \textit{Note:} A collaborative culture fosters open communication and teamwork, which encourages workers to embrace change as a collective effort rather than an individual challenge.
\item \textbf{Answer:} A
\\ \textit{Options:} A) Morale; B) Innovation; C) Growth resource; D) Adaptation
\\ \textit{Note:} Morale is not a key feature of the 'open systems' model of management because this model primarily focuses on the interactions and relationships between an organization and its external environment, rather than internal emotional or psychological factors.
\end{enumerate}

\section*{True / False}
\begin{enumerate}[label=\arabic*.]
\setcounter{enumi}{5}
\item \textbf{Answer:} False
\\ \textit{Options:} A) True; B) False
\\ \textit{Note:} Primary research refers to the collection of new data specifically for a particular research objective, rather than utilizing pre-existing data collected for different purposes.
\item \textbf{Answer:} True
\\ \textit{Options:} A) True; B) False
\\ \textit{Note:} Organisational structures characterized by democratic and inclusive styles of management are typically described as flat because they minimize hierarchical levels, promote collaborative decision-making, and empower employees, leading to a more egalitarian workplace.
\item \textbf{Answer:} False
\\ \textit{Options:} A) True; B) False
\\ \textit{Note:} The marketing period emphasizes understanding consumer needs and building relationships, which often includes addressing social concerns alongside maximizing sales and production efficiency.
\item \textbf{Answer:} False
\\ \textit{Options:} A) True; B) False
\\ \textit{Note:} Opinion polls are most useful in the formative and process stages of Macnamara's pyramid of evaluation, as they provide insights into public perception and feedback that can influence strategies before outcomes are assessed.
\item \textbf{Answer:} True
\\ \textit{Options:} A) True; B) False
\\ \textit{Note:} Exhibitions are designed specifically for sellers to display their offerings to potential buyers, facilitating direct interaction and promoting sales opportunities.
\end{enumerate}

\section*{Long Form Response}
\begin{enumerate}[label=\arabic*.]
\setcounter{enumi}{10}
\item \textbf{Answer:} In business administration, research design is typically categorized into three main types: exploratory, descriptive, and causal research. Exploratory research seeks to gain insights and understanding of a problem, descriptive research aims to describe characteristics of a population or phenomenon, while causal research investigates cause- and-effect relationships. Desk research, which involves the synthesis and analysis of existing information or data, is not traditionally classified among these categories because it does not involve the collection of new data through primary research methods. Instead, desk research serves as a supportive tool that can inform and guide the primary research process, making it more of a preliminary step rather than a standalone research design.
\\ \textit{Note:} correct because it accurately identifies the three main categories of research design in business administration-exploratory, descriptive, and causal-while clarifying that desk research, which relies on existing data rather than original data collection, does not fit within these traditional classifications.
\item \textbf{Answer:} Kantian ethics emphasizes the importance of treating individuals as ends in themselves, meaning that each person possesses intrinsic worth and should not be utilized merely as a means to achieve another's goals. In the context of business practices, this principle encourages ethical decision-making that respects the dignity and autonomy of all stakeholders, including employees, customers, and suppliers. By prioritizing the well- being and rights of individuals, businesses can foster trust, loyalty, and a positive reputation, ultimately leading to sustainable success. Therefore, adhering to this principle not only aligns with ethical imperatives but also contributes to long-term profitability and corporate responsibility.
\\ \textit{Note:} correct because it accurately captures Kantian ethics' central tenet of recognizing individuals' intrinsic worth, which mandates that ethical business practices must prioritize the dignity and autonomy of all stakeholders, thereby fostering a culture of trust and loyalty.
\end{enumerate}

\vfill
\noindent\textit{End of Answer Key}
\end{document}